So far we have seen codes of the type
\[
    \U^n \to \set{0,1}^*
\]

They are called Fixed-to-variable-length codes and all have error free recovery of the source from its representation.

We now discuss "Fixed-to-fixed" codes, we want
\[
    \C: \U^n \to \set{0,1}^k
\]

\begin{example}
    Suppose $\U=\set{a,b,c,d}$.
    \[
        \C: \U^n \to \set{0,1}^k
    \]
    $\card{\U^n} = 4^n$ and 
    $\card{\set{0,1}^k} = 2^k$
    
    error free recover requires $k\geqslant 2n$.
\end{example}

To obtain efficient codes, we will give up error-free recovery replace this by recovery with very small probability of error. For this, we will have the code assign binary representation only to a subset $\S \subseteq \U^n$. We will choose this subset to ensure 
\[
    \prob{(U_1\ldots U_n)\in\S} \simeq 1 \wedge \card{\S} \leqslant 2^k
\]

\begin{example}
    Suppose $\U = \set{a,b}p$, $p_a = 0.25$, $p_b = 0.75$. Source is memoryless stationary (iid) $U_1, U_2, \ldots$
    
    What kind of sequences $(u_1,\ldots,u_n)$ do we expect the source to produce?
    
    We expect that $(u_1,\ldots,u_n)$ to have $\simeq \frac{1}{4} n$ a's and $\simeq \frac{3}{4} n$ b's.
    
    Pick $\varepsilon > 0$ small. Let us define $\S = \set{\left.(u_1,\ldots,u_n) \middle\vert \length{(u_1,\ldots,u_n)}_a = \frac{1}{4}n(1\pm \varepsilon)\right.}$ .
    
    Intuitively, we expect $\set{(u_1,\ldots,u_n) \in \S}$ to be a high probability event.
    
    How about the size of $\S$?
    
    How many sequences $(u_1,\ldots,u_n)$ are there with $i$ a's and $n-i$ b's? 
    \[
        \binom{n}{i}
    \]
    So
    \[
        \begin{aligned}
            \card{\S} &= \sum\limits_{i = \ceil{\frac{n}{4}(1-\varepsilon)}}^{\floor{\frac{n}{4}(1+\varepsilon)}} \binom{n}{i}\\
            &\leqslant (n+1)\binom{n}{\floor{\frac{n}{4}(1+\varepsilon)}}
        \end{aligned}
    \]
    \[
        \begin{aligned}
            1 &= (p + (1-p))^n\\
            &=\sum\limits_{i=0}^n \binom{n}{i} p^i (1-p)^{n-i}\\
            &\geqslant \binom{n}{i} p^i (1-p)^{n-i}\\
            \Ra \binom{n}{i} &\leqslant \left(\frac{1}{p}\right)^i\left(\frac{1}{1-p}\right)^{n-i}\\
            &\leqslant \left(\frac{n}{i}\right)^i\left(\frac{n}{n-i}\right)^{n-i} \qquad \text{ with } p=\frac{i}{n}
        \end{aligned}
    \]
    
    \[
        \log\binom{n}{i} \leqslant n\left( \alpha \log\frac{1}{\alpha} + (1-\alpha)\log\frac{1}{1-\alpha} \right)
    \]
    \[
        \binom{n}{\floor{\frac{n}{4}(1+\varepsilon)}} \leqslant 2^{nH(U')}
    \]
    where $U' = \begin{cases}
        a & \text{ wp } \frac{1}{4}(1+\varepsilon)\\
        b & \text{ wp } 1 - \frac{1}{4}(1+\varepsilon)
    \end{cases}$.
    
    So $\card{\S} \leqslant (n+1)2^{nH(U')}$ so if we choose $k = aH(U') + \log(n+1)$ we will have $\card{} \leqslant 2^k$ with $\frac{k}{n} = \underbrace{H(U') + \frac{\log(n+1)}{n}}_{\simeq H(u) \text{if }\varepsilon\text{ is small and }n\text{ is large}}$ we will have "almost free" representation.
\end{example}

\subsection{General case}

We are given an alphabet $\U$ and a distribution $p$ on $\U$, $\set{p(u) \vert u\in\U}$, and we have a source that produces iid letters $U_1,U_2,\ldots$ each with distribution $p$. We want to have a set $T_{n,\varepsilon, p} \subseteq \U^n$ with the properties
\begin{enumerate}
    \item $\prob{(U_1\ldots U_n) \in T_{n,\varepsilon,p}} \simeq 1$
    \item $\card{T_{n,\varepsilon,p}} \leqslant ...$
\end{enumerate}

We say a sequence $(u_1\ldots u_n) \in \U^n$ to be $\varepsilon$-typical with respect to the distribution $p$ if for every $u\in\U, \frac{\card{\set{i \vert u_i = u}}}{n} = p(u) (1\pm \varepsilon)$.

Set $T_{n,\varepsilon, p}$ to be set of $(u_1,\ldots u_n)$'s which are $\varepsilon$-typical wrt $p$.

\begin{proposition}
    \[
        \prob{(U_1\ldots U_n) \not\in T_{n,\varepsilon,p}} \leqslant \frac{1}{n \varepsilon^2p_{\min}}
    \]
\end{proposition}
\begin{proof}
    Let $B_u$ be the vent that $\frac{1}{n}\length((U_1\ldots U_n))_u \not\in [p(u)(1-\varepsilon); p(u)(1+\varepsilon)]$. Then 
    \[
        \begin{aligned}
            \prob{(U_1\ldots U_n) \not\in T_{n,\varepsilon,p}} &= \prob{\bigcup\limits_{u\in\U} B_u}\\
            &\leqslant \sum\limits_{u\in\U} \prob{B_u}
        \end{aligned}
    \]
    
    $B_u$ is the event that $\left\lvert \frac{1}{n} \sum\limits_{i=1}^n \underbrace{\mathbb{1}\set{U_i = n}}_{X_i = \begin{cases}
    1 & \text{ wp } p(u)\\
    0 & \text{ wp } 1-p(u)
    \end{cases}} - p(u) \right\rvert > \varepsilon p(u)$.
    
    ie. $\left\lvert \frac{1}{n} \sum\limits_{i=1}^n X_i - p(u) \right\rvert > \varepsilon p(u)$.
    
    \[
        \prob{B_u} \leqslant \frac{1}{n\varepsilon^2 p_{\min}}
    \]
    where $p_{\min} = \min p(u)$.
\end{proof}

For the size of $T_{n, \varepsilon, p}$ observe first the followings
\begin{itemize}
    \item Fix a $(u_1,\ldots,u_n)$ in $T_{n,\varepsilon, p}$. Note that
    \[
        \begin{aligned}
            \prob{(u_1 \ldots u_n)=(U_1 \ldots U_n)} &= p(u_1)\cdots p(u_n)\\
            &= \prod\limits_{u\in\U}p(u)^{\overbrace{\length{u_1\ldots u_n}_u}^{np(n)(1\pm\varepsilon)}} \geqslant \prod\limits_u p(u)^{np(u)(1+\varepsilon)}\\
            &= \prod\limits_u 2^{np(u)\log p(u)(1+\varepsilon)} = 2^{-n(1+\varepsilon)H(u)}\\
            \Ra 1 &\geqslant \prob{(U_1\ldots U_n)\in T_{n,\varepsilon,p}}\\
            &= \sum\limits_{(u_1\ldots u_n \in T_{n,\varepsilon,p})} \prob{(U_1\ldots U_n) = (u_1\ldots u_n)}\\
            & \geqslant 2^{-n(1+\varepsilon)H(u)}\card{T_{n,\varepsilon,p}}\\
            \Ra \card{T_{n,\varepsilon,p}} &\leqslant 2 ^{(1+\varepsilon)nH(U)}
        \end{aligned}
    \]
\end{itemize}

\begin{theorem}
    If $U_1\ldots U_n$ are iid $\sim p$, then
    \[
        \forall \varepsilon > 0, \lim\limits_{n\to\infty}\prob{(U_1 \ldots U_n) \in T_{n, \varepsilon, p}} = 1
    \]
    
    Furthermore if $(u_1 \ldots u_n)\in T_{n,\varepsilon, p}$, then
    \[
        \begin{aligned}
            \prob{(U_1 \ldots U_n) = (u_1 \ldots u_n)} &= \prod\limits_{i=1}^n p(u_i)\\
            &= \prod\limits_{u\in\U} p(u) ^{\overbrace{\length{u_i}_u}^{np(n)(1\pm\varepsilon)}}\\
            &= 2^{-nH(u)(1\pm\varepsilon)}
        \end{aligned}
    \]
    ie.
    \[
        2^{-nH(u)(1+\varepsilon)} \leqslant \prob{(U_1 \ldots U_n) = (u_1 \ldots u_n)} \leqslant 2^{-nH(u)(1-\varepsilon)}
    \]
    Also $\card{T_{n,\varepsilon,p}} \leqslant 2^{nH(u)(1+\varepsilon)}$ and $\card{T_{n,\varepsilon,p}} \geqslant (1-\varepsilon) 2^{nH(u)(1-\varepsilon)}$
\end{theorem}

\begin{theorem}
    If $U_1\ldots U_n$ iid $\sim q$ and $(u_1\ldots u_n)\in T_{n,\varepsilon,p}$, then
    \[
        \begin{aligned}
            \prob{(U_1\ldots U_n) = (u_1\ldots u_n)} &= \prod\limits_{i=1}^n q(u_i)\\
            &= \prod\limits_{u\in\U} q(u)^{np(u)(1\pm \varepsilon)}\\
            &= 2^{-n(1\pm\varepsilon)\sum\limits_u p(u) \log\frac{1}{q(u)}}
        \end{aligned}
    \]
    But
    \[
        \begin{aligned}
            \sum\limits_u p(u) \log\frac{1}{q(u)} &= \sum\limits_u p(u) \log\frac{1}{p(u)} + \sum\limits_u p(u) \log\frac{p(u)}{q(u)}\\
            &= H(p) + \underbrace{D(p \Vert q)}_{\text{divergence}}
        \end{aligned}
    \]
    So
    \[
        \prob{(U_1\ldots U_n) = (u_1\ldots u_n)} = 2^{-(1\pm\varepsilon)n(D(q\Vert p + H(p))}
    \]
\end{theorem}

\begin{corollary}
    If $U_1\ldots U_n$ iid $\sim q$
    \[
        \begin{aligned}
            \prob{(U_1\ldots U_n)\in T_{n,\varepsilon,p}} &= \underbrace{2^{-(1\pm\varepsilon)n(H(p) + D(q\Vert p))}}_{\substack{\text{upper and lower bound on}\\\prob{(U_1\ldots U_n) = (u_1\ldots u_n)}}} \cdot \underbrace{(1\pm \varepsilon) 2 ^{nH(p)(1\pm \varepsilon)}}_{\substack{\text{upper and lower}\\\text{bound on }\card{T}}}\\
            &= 2^{-n(1\pm \varepsilon) D(p\Vert q)}2^{\pm n\varepsilon H(p)} (1\pm\varepsilon) \\
            &\doteq 2^{-nD(p\Vert q)}
        \end{aligned}
    \]
\end{corollary}


Recall also
\[
    \begin{aligned}
        I(X,Y) &= \sum\limits_{x,y} p(x,y) \log\frac{p(x,y)}{p(x)p(y)}\\
        &= D(p_{XY} \Vert p_X p_Y)
    \end{aligned}
\]

Consequently if $(X_1,Y_1)\ldots(X_n,Y_n)$ iid $p_Xp_Y$ then
\[
    \prob{((X_1,Y_1)\ldots(X_n,Y_n)) \in T_{n,\varepsilon, p_{XY}}} \doteq 2^{-n\overbrace{I(X;Y)}^{\substack{\text{computed}\\\text{wrt. } p_{XY}}}}
\]

\bigskip

"Application" of typical sequences in source coding. Given iid source $U_1U2\ldots$ $\sim p$. Pick $\varepsilon > 0$ (small) and $n$ (long) so that $T_{n,\varepsilon, p}$ has large probability ($>1-\varepsilon$). Pick $k = (1+\varepsilon)nH(u)$ so that $2^k$ is $\geqslant \card{T_{n,\varepsilon, p}}$. Consider the following source coding method.

First, assign to every member of $T_{n,\varepsilon,p}$ a distinct element of $\set{0,1}^k$, call this $\C : T_{n,\varepsilon,p} \to \set{0,1}^k$.

The source code is the following: if $(U_1\ldots U_n)\in T_{n,\varepsilon,p}$, represent it by $\C(U_1\ldots U_n)\in\set{0,1}^k$ else destroy the universe (it seems there is a nuclear button to do that).

This method uses $(1+\varepsilon) H(u) + \frac{1}{n}$ bits/letter and has a probability of decoding failure $< \varepsilon$.


\bigskip

"Dual" of Fix-to-variable source coding $\equiv$ variable to fixed length source coding (dictionary based source coding).

The idea: given an alphabet $\U$, we will find a dictionary $\D \subsetneq \U^*$, we will assign $\ceil{\log\card{\D}}$ bits binary representation to words in $\D$, wand then given $U_1U_2\ldots$ we will parse it into $w_1w_2\ldots$ of dictionary words, and represent each word by its binary description.

\begin{example}\label{codes_fixed_to_fixed:ex:1}
    $\U = \set{a,b,c}$, $\D = \set{\underbrace{a}_{000},\underbrace{ba}_{001},bb,bca,bcb,bcc,\underbrace{c}_{110}}$ requires no lookahead.
    
    \[
        bcaababccab\ldots \to \underbrace{bca}_{011} \underbrace{a}_{000} \underbrace{ba}_{001}\underbrace{bcc}_{101}\underbrace{a}_{000}
    \]
\end{example}

\begin{example}\label{codes_fixed_to_fixed:ex:2}
    $\U = \set{a,b,c}$, $\D=\set{aa,ab,b,c}$ $\ra$ invalid because it cannot parse $ac$.
\end{example}

\begin{example}\label{codes_fixed_to_fixed:ex:3}
    $\U = \set{a,b,c}$, $\D=\set{a,aa,aaa,b,c}$ requires lookahead.
    \[
        aabcab \to aa,b,c,a,b    
    \]
    first word as long as possible.
\end{example}

Tree representation of dictionaries.

Example \ref{codes_fixed_to_fixed:ex:1}:
\Tree [ [ $a$ ] [ [ $ab$ ][ $bb$ ][ [ $bca$ ][ $bcb$ ][ $bcc$ ] ] ] [ $c$ ] ]

Example \ref{codes_fixed_to_fixed:ex:2}:
\Tree [ [ [ $aa$ ] [ $ab$ ] ] [ $b$ ] [ $c$ ] ]

Example \ref{codes_fixed_to_fixed:ex:3}:
\Tree [ [ .$a$ [ .$aa$ [ .$aaa$ ] [ ] [ ] ] [ ] [ ] ] [] [ $c$ ] ]

The words of a valid \& no-lookahead (=prefix free) dictionaries form the leaves of a complete $\card{\U}$-ary tree.

Let us examine the working of a valid \& no-lookahead prefix-free dictionary.

\[
    \underbrace{U_1\ldots U_n}_{\text{iid.}} \to \underbrace{W_1}_{\in \D} \ldots W_m \ra \text{binary representation in }\ceil{\log\card{\D}}
\]

$W's$ are also iid.

\# of bits/letters $=\frac{m\log\ceil{\card{\D}}}{\sum\limits_{i=1}^n\length{W_i}} \to \frac{\ceil{\card{\D}}}{\esp{\length{W_1}}}$.

\begin{lemma}
    Suppose $Z$ is non negative integer valued RV, then
    \[
        \esp{Z} = \sum\limits_{n=0}^{\infty} \prob{Z > n}
    \]
\end{lemma}

Observation on the \# of word of a valid prefix free dictionary (=\# of leaves in the tree representation of the dictionary) and the \# of interior nodes of the tree.

\begin{tabular}{ccc}
    Tree & \# of interior nodes & \# of dictionary words\\
    \Tree [ [ $a$ ] [ $b$ ] [ $c$ ] ] & 1 & $\card{\U}$\\
    \Tree [ [ [ $aa$ ] [ $ab$ ] [ $ac$ ] ] [ $b$ ] [ $c$ ] ] & 2 & $2\card{\U} -1$\\
    \Tree [ [ $b$ ] [ [ $aa$ ] [ $ab$ ] [ [ $aca$ ] [ $acb$ ] [ $acc$ ] ] ] [ $c$ ] ] & 2 & $3\card{\U} - 2$
\end{tabular}

\# of words = $1 + (\card{\U} - 1) (\text{\# of interior nodes})$.



\begin{algorithm}
    \caption{\textsc{Tunstall}'s procedure to construct such a tree}
    \DontPrintSemicolon
    Start with $\D = \U$ with $\card{\U}$ words and 1 interior node (root)\;
    \While{$\card{\D} < M$}{
        convert the most probable word (leaf) into an interior node, grow $\U$ leaves from it.
    }
\end{algorithm}

\begin{figure}[!ht]
    \centering
    \includegraphics[scale=0.25]{figures/suchTree}
    \caption{Wow!}
\end{figure}

\FloatBarrier

\begin{lemma}
    Suppose $\D$ is valid prefix-free, suppose $U_1,U_2,\ldots$ are iid, then
    \[
        H(W_1) = \esp{\length{W_1}} H(U_1)
    \]
\end{lemma}
\begin{proof}
    \[
        \begin{aligned}
            H(W_1) &= \sum\limits_{w\in \D} P(w) \log \frac{1}{P(w)}\\
            &= \esp{\log\frac{1}{P(W_1)}}\\
            &= -\esp{\log P(W_1)}
        \end{aligned}            
    \]

    Also if $w=(w_1,\ldots,w_n)$ then
    \[
        \begin{aligned}
            P(w) &= \prod\limits_{i=1}^l P_v(u_i)\\
            \Ra P(W_1) &= \prod\limits_{i=1}^{\length{W_i}} P_v(U_i)\\
            \Ra H(W_1) &= \esp{\sum\limits_{i=1}^{l(U_1,U_2,\ldots)} - \log p_u(U_i)}\\
            &= \esp{\sum\limits_{i=0}^\infty (-\log p_u(U_i)) \mathbb{1}\set{i\leqslant l(U_1U_2\ldots)}}\\
            &= \sum_{i=0}^\infty \esp{ (-\log p_u(U_i)) \mathbb{1}\set{i\leqslant l(U_1U_2\ldots)}}
        \end{aligned}
    \]
\end{proof}

Recall now 
\[
    \begin{aligned}
        \text{\# bits/letter} &= \frac{ \ceil{\log\card{\D}}}{\esp{\length{W_1}}}\\
        &= \frac{\ceil{\card{\D}}}{H(W_1)} H(U_1)
    \end{aligned}
\]

To have any hope of showing that \# of bits/letter $\approx H(U_1)$ we need to show that $H(W_1) \geqslant \log \card{\D} - \text{something we can control}$.

To show this, observe how the \textsc{Tunstall}'s procedure picks the interior nodes. Initially we have one interior node of weight 1 (root). Let $Q_1,Q_2,\ldots$ the weight of the interior nodes with $1=Q_1 \geqslant Q_2 \geqslant \ldots \geqslant Q_k$ $\ra$ last intermediate work picked. At the point that $k$-th intermediate word is picked, all words have $P(\text{word}) \leqslant Q_k$ the statement remains true after one pick. Also for any word $P(\text{word}) = Q_i P(\text{last letter of word}) \geqslant Q_k \min\limits_up_u(u)$.

\[
    2p_0 \leqslant \prob{W=w} \leqslant q
\]
where
$p_0= \min\limits_x p_u(x)$ and $q$ the weight of the last picked interior node.

\begin{lemma}
    If $W$ is a random variable taking values in $\D$ and if for each $w\in\D$
    \[
        qp_0 \leqslant \prob{W=w}\leqslant q
    \]
    then
    \[
        \log \card{\D}  - \log \frac{1}{p_0} \leqslant H(W) \leqslant \log \card{\D}
    \]
\end{lemma}
\begin{proof}
    \[
        \begin{aligned}
            H(W) &= \sum\limits_w p(w) \log \frac{1}{p(w)}\\
            &\geqslant \sum_w p(w) \log \frac{1}{q}\\
            &< \log \frac{1}{q}
        \end{aligned}
    \]
    Note also
    \[
        \begin{aligned}           
            qp_0 \leqslant \prob{W=w} &\Ra qp_0\card{\D}\sum\limits_{w\in\D} qp_0 \leqslant \sum\limits_{w \in D}\prob{W=w} = 1\\
            &\Ra q\leqslant \frac{1}{p_0\card{D}}\\
            &\Ra \log \frac{1}{q} \geqslant \log\card{\D} + \log p_0
        \end{aligned}
    \]
\end{proof}

Using the lemma, we have $H(W) \geqslant \log\card{D} - \log \frac{1}{p_0}$. So, for the $\D$ constructed by \textsc{Tunstall}
\[
    \text{\# of bits/lettre} = H(U) \frac{\ceil{\log\card{\D}}}{H(W)} \leqslant H(W) \frac{\log\card{\D}+1}{\log\card\D - \log\frac{1}{p_0}} \underset{\card\D \to \infty}{\longra} 1
\]

Consequently, we have
\begin{corollary}
    Given $\U$, and $p_u$, $\varepsilon > 0$, there exists a dictionary based code for which \# bits/letter $\leqslant H(u)(1+\varepsilon)$.
\end{corollary}